%% kapa: exam=1 round=27 year=2016 period=2 form=A subject=accounting subject_ko=회계학 seq=9 number=49
%% source: Q-Net 감정평가사 기출문제 (cst00309) artlSeq=5169402 file=extracted/1008/5169402/제27회 감정평가사 1차 문제/제27회 감정평가사 2교시문제(A형).pdf
%% conversion: kordoc 4.24.51 -> tools/gen1.py (auto); review=proofread-sonnet; flags=none
\begin{question}{49}
(주)감평은 기계장치를 (주)대한의 기계장치와 교환하였다. 교환시점에 두 회사가 소유하고 있던 기계장치의 장부금액과 공정가치는 다음과 같다.

{\small\centering
\begin{tabular}{|c|r|r|}\hline
구 분 & \multicolumn{1}{c|}{(주)감평} & \multicolumn{1}{c|}{(주)대한} \\ \hline
\multicolumn{1}{|l|}{취득원가} & ￦1,000,000 & ￦1,200,000 \\ \hline
\multicolumn{1}{|l|}{감가상각누계액} & 300,000 & 600,000 \\ \hline
\multicolumn{1}{|l|}{공정가치} & 600,000 & - \\ \hline
\end{tabular}\par}

이 기계장치의 교환과 관련하여 (주)감평은 (주)대한으로부터 현금 ￦50,000을 추가로 수령하였다. (주)감평이 교환거래로 인식해야할 처분손익은? (단, 교환거래는 상업적 실질이 있다.)
\begin{choices}
  \item 처분이익 ￦50,000
  \item 처분손실 ￦50,000
  \item 처분이익 ￦100,000
  \item 처분손실 ￦100,000
  \item 처분손실 ￦150,000
\end{choices}
\end{question}
\answer{4}
